\section{Negative results}\label{sec:negative}

\para{Ternary weights that fit in L3 are not faster at batch.} We trained the multirate model with
ternary weights $\{-1,0,+1\}$ and per-channel absmean scales by quantisation-aware fine-tuning
from a 100\,h checkpoint~\cite{ma2024bitnet}, with activations fake-quantised exactly as the
engine quantises them. The ternary model costs \TernaryWerCost\pp WER and stores its encoder in
\TernaryMiB\,MiB, which fits a Zen~2 CCX's 16\,MiB L3. The engine's ternary lookup kernel wins at a
single stream, where weight traffic dominates. At batch it loses: \TernaryDenCCX{} streams per CCX
against \IntEightMRDenCCX{} for INT8 multirate on x86, both measured with the engine's development
tool (not the evaluation harness, so not comparable with \cref{tab:density}). Once sixteen or more streams share each weight traversal
the engine is compute-bound, weight bytes stop mattering, and the int8 dot-product units do more
useful work per cycle than table lookups.

\para{Removing instructions does not remove time.} On AVX2 the signed int8 GEMM spends three of
ten instructions per 64 multiply-accumulates on taking absolute values and transferring signs. An
unsigned-offset 7-bit tile removes all three and is bit-exact against our 7-bit reference. On Zen~2 it is
\emph{slower}: \UsevenVsKzeroLo--\UsevenVsKzeroHi{}$\times$ the speed of the signed tile across
layer shapes and hosts, against a pre-registered bar of \UsevenBar$\times$. It keeps four vector
multiplies per 64 MACs; the 7-bit activation variant, which drops one multiply, is the only one that
wins, by about \ActSevenGain\%. On this core time per MAC follows the multiply count, not the
instruction count.

\para{Bit-identical non-GEMM cuts are small.} Hoisting rotary tables and fusing residual adds
with the next normalisation are bit-identical and kept, but gain only \NonGemmGain\% end to end,
inside run-to-run noise: the remaining non-GEMM time is real attention work, not overhead.
